\documentclass[12pt]{article} % Default font size is 12pt, it can be changed here

\usepackage[spanish]{babel} % Configura el idioma del documento a español

\usepackage[a4paper, margin=1in]{geometry}

\usepackage{graphicx} % Required for including pictures

\usepackage{url}

\usepackage{float} % Allows putting an [H] in \begin{figure} to specify the exact location of the figure
\usepackage{wrapfig} % Allows in-line images such as the example fish picture

\linespread{1.2} % Line spacing

%\setlength\parindent{0pt} % Uncomment to remove all indentation from paragraphs

\graphicspath{{./imagenes/}} % Specifies the directory where pictures are stored

%para escribir en español
\usepackage[utf8x]{inputenc}

\usepackage{amsmath}
\usepackage{multirow}
\usepackage{longtable} % Para tablas que pueden ocupar más de una página
\usepackage{array} % Para opciones avanzadas de tablas
\usepackage{booktabs} % Para mejorar la calidad visual de las tablas
\usepackage{listings}
\usepackage{fancyvrb}
\DefineVerbatimEnvironment{Verbatim}{Verbatim}
{breaklines=true} % Activa el ajuste automático de líneas
\usepackage{caption}

%opening

\begin{document}
	\begin{titlepage}
		\newcommand{\HRule}{\rule{\linewidth}{0.5mm}} % Defines a new command for the horizontal lines, change thickness here
		
		\center % Center everything on the page
		
		\textsc{\LARGE Programación Lógica}\\[1.5cm] % Name of your university/college
		\textsc{\Large Curso 2025}\\[0.5cm] % Major heading such as course name
		\textsc{\large Grupo 01}\\[0.5cm] % Minor heading such as course title
		
		\HRule \\[0.4cm]
		{\huge \bfseries Informe - Laboratorio 2}\\[0.4cm] % Title of your document
		\HRule \\[1cm]
		
		\begin{minipage}{0.4\textwidth}
			\begin{flushleft} \large
				\emph{Autores:}\\
				\textsc{Gonzalo Apkarian Alonso \\ 5.494.604-2}\\ % Your name
				\textsc{Diego Pisa Sánchez \\ 5.510.635-2}\\ % Your name
				\textsc{Leonardo Pesce López \\ 5.471.535-2}\\
				\textsc{Mauricio Morales Gonzalez \\ 5.278.642-4}
			\end{flushleft}
		\end{minipage}
		~
		\begin{minipage}{0.4\textwidth}
			\begin{flushright} \large
				\emph{Profesores:} \\ 
				Luis \textsc{Chiruzzo} \\ 
				Aiala \textsc{Rosá} \\ 
				Juan Pablo \textsc{Conde}
			\end{flushright}
		\end{minipage}\\[3cm]
		
		% Logo
		\begin{figure}[H]
			\centering
			\includegraphics[width=0.2\textwidth]{logo.png} % Ajusta el tamaño aquí (0.2\textwidth es un ejemplo)
			\label{fig:logo}
		\end{figure}
		
		% Fecha
		{\large \today}\\[1cm] % Puedes cambiar \today por una fecha específica si lo prefieres
		
		\vfill % Esto asegura que el contenido se distribuya adecuadamente en la página
		
	\end{titlepage}
	
	\section{Módulos}
	Para este laboratorio no se usaron módulos extra, se implementaron todos los predicados en el mismo archivo principal. Esta decisión se tomó debido a que se tornó una complicación decidir como modularizar de manera eficiente el programa y no se contaba con mucho tiempo, por lo que el equipo se fue por el camino más rápido aunque sin seguir patrones de un buen diseño en su totalidad.
	
	Sin embargo, para facilitar la lectura y corrección del programa, cada vez que se iniciaba el desarrollo de un predicado principal se dejaban líneas con caracteres comentados de forma que se pueda diferenciar fácilmente en el IDE las distintas funciones.
	
	\section{Predicados principales}
	En esta sección se mencionarán y explicarán, además del funcionamiento de los predicados necesarios para que funcione \textit{ceritolog} (es decir, los que vienen explicados en la letra), aquellos predicados auxiliares que el equipo considera esenciales para que el principal pueda realizar su función.
	
		\subsection{Tablero}
		En este predicado simplemente se pide devolver un tablero de tamaño N vacío, donde N es la cantidad de vértices que tiene el tablero por lado, generando una matriz cuadrada donde en realidad quedan $(N-1)\cdot(N-1)$ celdas. \\
		
			\subsubsection{Diseño del tablero}
			El tablero se crea usando las celdas, la cual es una tupla de 3 elementos (arista superior, arista izquierda, jugador) o (F, C, J). Aquí se tuvo que hacer la primera decisión del diseño del tablero, ya que, si bien el tablero es $N\cdot N$ en vértices, se tomó la decisión de crearlo con $N\cdot N$ celdas para poder tener la información de las aristas de los bordes inferior y derecho. Para esas celdas extra, en la última fila solo se tiene acceso a modificación de la arista superior (que sería la arista inferior del tablero) mientras que en la última columna solo se puede modificar la arista izquierda de esas celdas (que sería la arista derecha del tablero).
			
		\subsection{Fin del juego}
		En este predicado se retorna el final del juego, cuál jugador ganó y con cuantos puntos terminó cada uno. En caso de que no sea el final del juego, el predicado falla.
		
			\subsubsection{Diseño de "Fin del juego"}
			Para el fin del juego se tomó que todas las líneas deben estar marcadas, es decir, una vez que se inicia una partida se debe terminar por completo. Para ello se utiliza un predicado auxiliar que justamente se encarga de verificar que estén todas las líneas marcadas.\\
			
			\subsubsection{Predicado auxiliar 1 - todas\_lineas\_marcadas}
			El predicado auxiliar \textit{todas\_lineas\_marcadas} se encarga de verificar que todas las celdas internas (las que estan entre $1$ y $(N-1)$) estén en marcadas como $c(1,1,\_)$, mientras que las de la última fila $c(1,0,0)$ y las de la última columna $c(0,1,0)$. Es importante que las celdas de los bordes derecho e inferior no tengan un jugador que las haya conquistado, de otra manera no estaría funcionando bien el juego.\\
			
			\subsubsection{Predicado auxiliar 2 - contar\_puntos}
			El otro predicado auxiliar que usa esta función es el de \textit{contar\_puntos} que pasa celda por celda sumando en el contador del jugador correspondiente si está conquistada por él. Por ejemplo, si se tiene la celda $c(\_,\_,1)$, se suma un punto al jugador $1$. Para que esto funcione se debe verificar antes que el juego está terminado de forma que todas las aristas estén marcadas debido a que $c(1,0,1)$ no es una instancia válida de una casilla, así como tampoco lo es $c(0,1,1)$ ni $c(0,0,1)$ (en todas se usa como ejemplo el jugador $1$, pero si en lugar de $1$ fuera $2$ tampoco serían viables estas instancias).
			
		
	
\end{document}